\chapter{Aparência: Mudando o Visual}

\begin{objetivos}
\begin{itemize}
\item Fazer o ator falar e pensar com balões;
\item Trocar de fantasia para criar movimento;
\item Mudar tamanho, cor e esconder ou mostrar o ator.
\end{itemize}
\end{objetivos}

\section{Falar e pensar}

\begin{pensaassim}[Pense assim: uma história em quadrinhos]
Nos quadrinhos, um balão mostra o que o personagem diz e outro mostra o que ele pensa. No Scratch, você também escolhe o texto e o tipo de balão. Esses balões aparecem na tela: eles não produzem voz.
\end{pensaassim}

Os blocos roxos \bloco[scLooks]{Aparência} cuidam do \textbf{visual} do ator. Os dois mais simples fazem aparecer um balão de fala ou de pensamento:

\begin{center}
\renewcommand{\arraystretch}{1.8}
\begin{tabular}{@{}ll@{}}
\bloco[scLooks]{diga \entrada{Olá!} por \entrada{2} segundos} & fala e depois o balão some \\
\bloco[scLooks]{diga \entrada{Olá!}} & fala e o balão fica na tela \\
\bloco[scLooks]{pense \entrada{Hmm...} por \entrada{2} segundos} & balão de pensamento \\
\end{tabular}
\end{center}

\begin{dica}
Use \bloco[scLooks]{diga \entrada{}} com o texto vazio para apagar um balão que ficou na tela.
\end{dica}

\section{Fantasias}

Cada ator pode ter várias \textbf{fantasias}. Uma fantasia é uma imagem diferente do mesmo ator, como os quadros de um desenho animado. O gatinho já vem com duas: uma de perna aberta e outra de perna fechada.

\begin{vocabulario}[Vocabulário: fantasia]
\textbf{Fantasia} é cada aparência de um ator. Trocar de fantasia depressa dá a sensação de movimento, e isso se chama \textbf{animação}.
\end{vocabulario}

\begin{pensaassim}[Pense assim: trocando de roupa]
Você pode vestir um pijama ou uma roupa de passeio e continuar sendo a mesma pessoa. As fantasias são aparências diferentes do mesmo ator: trocar a fantasia não cria outro personagem.
\end{pensaassim}

\begin{center}
\renewcommand{\arraystretch}{1.8}
\begin{tabular}{@{}ll@{}}
\bloco[scLooks]{próxima fantasia} & passa para a fantasia seguinte \\
\bloco[scLooks]{mude para a fantasia \menu{gato-b}} & escolhe uma fantasia específica \\
\end{tabular}
\end{center}

\begin{missao}[Missão 3: O gatinho que corre]
Faça o gatinho correr pelo palco, alternando as fantasias a cada passo.
\end{missao}

\begin{codigo}[Gatinho andando]
\begin{pilha}
\bk[scEvents]{quando \bandeira\ for clicado}
\sbc[scControl]{repita \entrada{10} vezes}
\bk[scMotion]{mova \entrada{10} passos}
\bk[scLooks]{próxima fantasia}
\bk[scControl]{espere \entrada{0.1} segundos}
\sbcend
\end{pilha}
\end{codigo}

\begin{cuidado}
Sem o bloco \bloco[scControl]{espere \entrada{0.1} segundos}, as fantasias trocam rápido demais e o gatinho parece tremer em vez de correr.
\end{cuidado}

\section{Tamanho, cor e visibilidade}

\begin{pensaassim}[Pense assim: uma foto e um fantoche]
Ao ampliar uma foto ou aplicar um filtro, você muda o tamanho ou as cores da imagem. Os efeitos do ator fazem algo parecido. Já esconder um ator é como colocar um fantoche atrás da cortina: você deixa de vê-lo, mas ele continua no projeto e pode aparecer de novo.
\end{pensaassim}

\begin{center}
\renewcommand{\arraystretch}{1.8}
\begin{tabular}{@{}ll@{}}
\bloco[scLooks]{mude o tamanho para \entrada{150} \%} & define o tamanho exato (100\% é o normal) \\
\bloco[scLooks]{mude o tamanho em \entrada{10}} & cresce (ou encolhe, se for negativo) \\
\bloco[scLooks]{mude o efeito \menu{cor} para \entrada{25}} & troca a cor do ator \\
\bloco[scLooks]{esconda} & o ator some do palco \\
\bloco[scLooks]{mostre} & o ator aparece de novo \\
\end{tabular}
\end{center}

\begin{curiosidade}
Efeitos como \textbf{cor}, \textbf{olho de peixe} e \textbf{mosaico} deixam o ator com cara de filtro de foto! Para voltar ao normal, use o bloco \bloco[scLooks]{tire os efeitos gráficos}.
\end{curiosidade}

\section{Hora de praticar}

\begin{exercicio}[Pensando alto]
Faça o gatinho \textbf{pensar} ``Que dia lindo!'' por 3 segundos e depois \textbf{dizer} ``Vamos brincar!'' por 2 segundos.
\end{exercicio}

\begin{solucao}
\begin{pilha}
\bk[scEvents]{quando \bandeira\ for clicado}
\bk[scLooks]{pense \entrada{Que dia lindo!} por \entrada{3} segundos}
\bk[scLooks]{diga \entrada{Vamos brincar!} por \entrada{2} segundos}
\end{pilha}
\end{solucao}

\begin{exercicio}[Gigante e pequeno]
Faça o gatinho ficar com o dobro do tamanho (200\%), esperar 1 segundo e depois voltar ao tamanho normal (100\%).
\end{exercicio}

\begin{solucao}
\begin{pilha}
\bk[scEvents]{quando \bandeira\ for clicado}
\bk[scLooks]{mude o tamanho para \entrada{200} \%}
\bk[scControl]{espere \entrada{1} segundos}
\bk[scLooks]{mude o tamanho para \entrada{100} \%}
\end{pilha}
\end{solucao}

\begin{exercicio}[Pisca-pisca]
Faça o gatinho \textbf{sumir e aparecer} 5 vezes, esperando meio segundo entre cada mudança.
\end{exercicio}

\begin{solucao}
\begin{pilha}
\bk[scEvents]{quando \bandeira\ for clicado}
\sbc[scControl]{repita \entrada{5} vezes}
\bk[scLooks]{esconda}
\bk[scControl]{espere \entrada{0.5} segundos}
\bk[scLooks]{mostre}
\bk[scControl]{espere \entrada{0.5} segundos}
\sbcend
\end{pilha}
\end{solucao}

\begin{desafio}[Gato arco-íris]
Faça o gatinho mudar de cor 12 vezes seguidas, somando 10 ao efeito de cor a cada vez, com uma pequena pausa entre as trocas. Dica: use \bloco[scLooks]{mude o efeito \menu{cor} em \entrada{10}}.
\end{desafio}

\begin{solucao}
\begin{pilha}
\bk[scEvents]{quando \bandeira\ for clicado}
\sbc[scControl]{repita \entrada{12} vezes}
\bk[scLooks]{mude o efeito \menu{cor} em \entrada{10}}
\bk[scControl]{espere \entrada{0.2} segundos}
\sbcend
\end{pilha}
\end{solucao}
